\section{A conditional three-parameter section}
The full V20B/M3 equations do not admit an exact universal reduction to
three independent parameters: already $U_1,U_2,V_2,\Omega_2$ supply four
independent experimental controls at fixed $R,\Omega_1$ and particle data.
They multiply two distinct RF families, a static DC force, and the second
external phase rate. Their effects are functionally independent on generic
states.

For a clearly specified section, fix $\nu$, the MFS spectrum and geometry,
$s,Q,C_{12},\lambda_1,\lambda_2,h_1,h_2,g_z$, particle radii, and RF phases
as initial-state variables. Let only $A,B,\beta$ vary. A bijective,
dimensionless parameterization of this section is
\begin{equation}
\Pi_T=\frac{B^2}{A^2},\qquad
\Pi_D=\frac{A^2}{\lambda_1},\qquad
\Pi_R=\frac{\beta}{A^2},\qquad A,B,\lambda_1>0.
\end{equation}
The inverse is
\begin{equation}
A=\sqrt{\Pi_D\lambda_1},\qquad
B=\sqrt{\Pi_T\Pi_D\lambda_1},\qquad
\beta=\Pi_R\Pi_D\lambda_1.
\end{equation}
$\Pi_T$ is a nominal squared drive ratio, not the actual RF torque ratio;
$\Pi_D$ compares the first RF drive to linear drag; $\Pi_R$ compares the
DC coefficient with the first drive squared. The three coordinates retain
all dependence on $(A,B,\beta)$ within the fixed section. They do not
remove $\nu$, particle contrasts, Coulomb, nonlinear drag, or geometry.
Any boundary inferred from this section must name all fixed groups and the
model used to obtain it.
